\documentclass[10pt,a4paper]{amsart}
\usepackage[margin=19mm]{geometry}
\usepackage{amsmath,amssymb,mathtools}
\usepackage[scr=rsfs,cal=euler]{mathalfa}
\usepackage{xfrac}
\usepackage[unicode,hidelinks,bookmarks=false]{hyperref}
\newcommand{\ie}{i.e.\ }
\newcommand{\cf}{cf.\ }
\newcommand{\resp}{respectively\ }
\newcommand{\Subsection}{\subsection}
\providecommand{\nobreakdash}{\nobreak\hbox{-}}
\setlength{\emergencystretch}{2em}
\multlinegap0pt
\usepackage{amssymb}
\usepackage{xcolor}
\usepackage{enumerate}
\usepackage{float}
\newtheorem{proposition}{Proposition}
\newtheorem{lemma}{Lemma}
\title{Differentiable normal linearization of partially hyperbolic dynamical systems}
\author{Weijie Lu, Yonghui Xia, Weinian Zhang, Wenmeng Zhang}
\date{Source: arXiv:2603.07063v1 (CC BY 4.0)}
\begin{document}
\maketitle

\noindent\emph{Source and licence.} Differentiable normal linearization of partially hyperbolic dynamical systems. Version arXiv:2603.07063v1 (CC BY 4.0), distributed by arXiv under the Creative Commons Attribution 4.0 International licence, http://creativecommons.org/licenses/by/4.0/, as recorded in the arXiv licence metadata for this exact version and shown on the abstract page https://arxiv.org/abs/2603.07063v1. Full licence notice: \texttt{licenses/arxiv-2603.07063-CC-BY-4.0.txt}.\par\smallskip

\noindent\textit{Editorial scope.} Counted unit: the article's lemma dec-lem (source0.tex lines 995--1018). The article's complete original proof of it is delivered with it (source0.tex lines 1027--1095, one \texttt{proof} environment of 69 lines, which the article places away from the statement). No mathematics is written, repaired or re-derived: the statement and the proof are the article's own lines, in the article's own notation, and no retained line is substituted or re-flowed.  Retained uncounted as the source-local context the proof needs: lines 451--454; lines 514--520; lines 530--536; lines 615--629; lines 991--993.  Nothing was omitted from any retained span, so the package carries no omission record.  No retained line contains a citation, so the wrapper carries no bibliography and no bibliography entry.  No retained line contains a figure environment, an image-inclusion command or drawing code, so no image is delivered and the package declares no asset dependency.  Editorial formatting only, in addition: an AMS \texttt{amsart} preamble that restates the article's own declarations and macros used by the retained lines, namely the environments \texttt{proposition}, \texttt{lemma} on separate counters, together with the standard packages those lines need.  The retained environments print in this document as Proposition 1, Lemma 1 in order of appearance, whereas the article prints the same items with the numbering of its own full document.  The complete native source of the article as distributed in the arXiv e-print for this exact version is bundled unchanged as \texttt{sources/195915/source0.tex}.
\begin{align}\label{Df-hold}
 \|Df(x)\| \le \delta_f, \quad   
 \|D  f(x)-D  f(y)\| \le M \|x-y\|^\alpha,\quad \forall x,y\in\mathbb{R}^d,
\end{align}

\begin{align}\label{ED}
\begin{split}
\|\mathcal{A}(m,n;x_c)\pi_s\| \le K (\lambda_s^+)^{m-n}, \quad & \forall m\ge n,
\\
\|\mathcal{A}(m,n;x_c)\pi_u\| \le K (\lambda_u^-)^{m-n}, \quad & \forall m\le n,
\end{split}
\end{align}

\begin{align}\label{g-center}
\begin{split}
 \|\mathcal{A}(m,n;x_c)\pi_c\| \le K (1+\varsigma)^{m-n}, \quad & \forall m\ge n,
\\
\|\mathcal{A}(m,n;x_c)\pi_c\| \le K (1-\varsigma)^{m-n}, \quad & \forall m\le n.
\end{split}
\end{align}

\begin{proposition}\label{Fol-thm}
Suppose that $F$ is a $C^{1,\alpha}$ diffeomorphism
 such that \eqref{ED}-\eqref{g-center} hold.
Then there exists an unstable foliation of $F$ with leaves
\begin{align}\label{unst-foli}
\mathcal{W}_u(x)=\{z_u+h_u(x,z_u): z_u\in X_u \},\quad \forall x\in\mathbb{R}^d,
\end{align}
where $h_u: \mathbb{R}^d \times X_u \to X_{cs}$ is continuous in both variables and is $C^{1,\alpha}$-smooth in $z_u$ with a globally bounded derivative.
Furthermore, $h_u$ satisfies that $h_u(x,\pi_u x)=\pi_{cs} x$ and
\begin{align}\label{hh-u}
\|h_u(x,z_u)-\pi_{cs}x\|
\le L\|(\pi_{su} x,z_u)\|^{1+\beta},\quad \forall x\in \mathbb{R}^d,~\forall z_u\in X_u,
\end{align}
where $0< \beta<(\log(1+\varsigma)-\log\lambda_u^-)/\log\lambda_s^-$.
\end{proposition}

\begin{align}\label{Dec-1}
H(x):=x_u+ h_u(x_{cs}, x_u), \quad \forall x=(x_{cs},x_u)\in\mathbb{R}^d,
\end{align}

\begin{lemma}\label{dec-lem}
Suppose that all conditions of Proposition {\rm \ref{Fol-thm}} hold. Then, the inverse of $H$ exists such that
\begin{align}\label{Dec-2}
H^{-1}(x)=x_u+h_u(x,0), \quad \forall x\in\mathbb{R}^d,
\end{align}
both of which satisfy that $H^{\pm 1}(x_{cs})=x_{cs}$ and
\begin{align}\label{Dif-3}
\|H^{\pm 1}(x)-x\|\le L\|x-\tilde{x}_c\|^{1+\beta},\quad \forall \tilde{x}_c\in X_c.
\end{align}
Moreover,
$
\hat{F}(x):=H^{-1} \circ F \circ H(x)
$
is a continuous mapping such that
\begin{align}\label{xu-F}
\begin{split}
&\pi_{cs} \hat{F}(x)=F(x_{cs}),\quad \|\partial_{x_u}(\pi_u\hat{F})(x)-A_u\|\le \delta,
\\
&\|\partial_{x_u}(\pi_u\hat{F})(x_{cs},x_u)-\partial_{x_u}(\pi_u\hat{F})(x_{cs},\tilde x_u)\|\le M\|x_u-\tilde x_u\|^\beta %\quad \forall x_{cs}\in X_{cs},~x_u,\tilde x_u\in X_u,
\end{split}
\end{align}
for all $x_{cs}\in X_{cs}$ and $x_u,\tilde x_u\in X_u$,
where $\delta>0$ is a small constant and $M>0$ is a constant.
\end{lemma}

\begin{proof}[Proof of Lemma {\rm \ref{dec-lem}.}]
In order to show that
$
H^{-1}\circ H=H \circ H^{-1}=id,
$
we note from \eqref{Dec-1} that
$\pi_uH(x)=x_u$ and $H(x)\in \mathcal{W}_u(x_{cs})$, implying that $H(x)$ and $x_{cs}$ belong to the same leaf, i.e.,
\begin{align}\label{huhu}
h_u(H(x),z_u)=h_u(x_{cs},z_u),
\end{align}
and note from Proposition \ref{Fol-thm} that $h_u(x,x_u)=x_{cs}$.
Then, by \eqref{Dec-2} we have
\begin{align*}
H^{-1}(H(x)) =\pi_u H(x) + h_u(H(x),0)
=x_u+h_u(x_{cs},0)=x_u+x_{cs}=x.
\end{align*}
Moreover, noting that $\pi_{cs}H^{-1}(x)=h_u(x,0)$ and $x=x_u+h_u(x,x_u)$ belong to the same leaf, we have
\begin{align*}
H(H^{-1}(x))=x_u+h_u(\pi_{cs}H^{-1}(x),x_u)=x_u+h_u(x,x_u)=x.
\end{align*}
Therefore, $H:\mathbb{R}^d\to \mathbb{R}^d$ is a homeomorphism.

In order to prove \eqref{Dif-3}, by \eqref{Dec-1}-\eqref{Dec-2}, we see that $H^{\pm 1}(x_{cs})=h_u(x_{cs},0)=x_{cs}$.
Moreover, it follows from \eqref{Dec-1} and \eqref{hh-u} (with $x=x_{cs}$) that
\begin{align}\label{Dec-10}
\frac{\|H(x)-x\|}{\|x-x_c\|^{1+\beta}}
=&~  \frac{\|x_u+h_u(x_{cs},x_u)-x\|}{\|x_{su}\|^{1+\beta}}
\notag \\
= &~ \frac{\| h_u(x_{cs},x_u)-x_{cs}\|}{\|(x_s,x_u)\|^{1+\beta}}\le L\quad {\rm for}~x\ne x_c.
\end{align}
Then, for any $\tilde{x}_c\in X_c$, we see from \eqref{Dec-10} that
\begin{align}\label{Dec-11}
\frac{\|H(x)- x\|}{\|x-\tilde{x}_c\|^{1+\beta}}
\le \frac{\|H(x)- x\|}{\|x-x_c\|^{1+\beta}}\le L\quad {\rm for}~x\ne \tilde x_c~{\rm and}~x\ne x_c
\end{align}
since $\|x-x_c\|=\|x_{su}\| \le \|(x_{su}, x_c-\tilde{x}_c)\|= \|x-\tilde{x}_c\|$.
Similarly, for $H^{-1}$, we see from \eqref{Dec-2} and \eqref{hh-u} (with $z_u=0$) that
\begin{align*}
\frac{\|H^{-1}(x)-x\|}{\|x-x_c\|^{1+\beta}}
=&~  \frac{\|x_u+h_u(x,0)-x\|}{\|x_{su}\|^{1+\beta}}
\\
=&~  \frac{\| h_u(x,0)-x_{cs}\|}{\|(x_{su},0)\|^{1+\beta}}\le L\quad {\rm for}~x\ne x_c,
\end{align*}
and then \eqref{Dif-3} holds.

Moreover, by the invariance of the unstable foliation and by \eqref{huhu}, we know that $F\circ H(x)$
and $F(x_{cs})$ belong to the same leaf, i.e.,
$
h_u(F\circ H(x),z_u)=h_u(F(x_{cs}),z_u).
$
Thus,
\begin{align*}
\pi_{cs}\hat F(x)=&~\pi_{cs}H^{-1} \circ F \circ H(x) =h_u(F \circ H(x),0)
\\
=&~ h_u(F(x_{cs}),0)=F(x_{cs}).
\end{align*}
Moreover, since \eqref{Dec-1}-\eqref{Dec-2} give
$
\pi_{u}\hat F(x)=\pi_{u}F \circ H(x)=\pi_{u}F(x_u+ h_u(x_{cs}, x_u)),
$
we compute that
\begin{align*}
\partial_{x_u} (\pi_{u}\hat F)(x)=&~ \Big(\partial_{x_{cs}}(\pi_{u}F)(x_u+ h_u(x_{cs}, x_u))\partial_{x_u} h_u(x_{cs}, x_u),
\\
&\quad  \partial_{x_u}(\pi_{u}F)(x_u+ h_u(x_{cs}, x_u))\Big),
\end{align*}
which together with \eqref{Df-hold} implies \eqref{xu-F}.
This completes the proof.
\end{proof}

\end{document}
